% =====================================================================
%  4.1  Atividade 3.2.2.1 -- Implantar modelos de gestão
% =====================================================================
\subsection{Atividade 3.2.2.1 -- Implantar modelos de gestão}
\label{subsec:modelos-gestao}

\subsubsection{Objetivo e fundamentação}

Definir e implantar os arranjos de gestão que garantirão a sustentabilidade
operacional e institucional dos sistemas de abastecimento de água e de
esgotamento sanitário implantados no âmbito do PSH-PR.

A dispersão das comunidades rurais e as limitações de capacidade operacional
dos municípios exigem arranjos institucionais que promovam cooperação entre
os diferentes níveis de governança. Os modelos de gestão específicos para o
saneamento rural encontram-se em definição. Como alternativa preliminar,
analisa-se um modelo baseado em arranjos intermunicipais, possivelmente
estruturado por meio de consórcios públicos ou mecanismos institucionais
equivalentes, com o objetivo de fortalecer a capacidade de gestão dos
municípios e garantir suporte técnico às comunidades \cite{mop2026}.

Prevê-se a atuação integrada do IAT, da SECID e do IDR-Paraná no
planejamento, na coordenação técnica, na mobilização social e no
monitoramento; os municípios participam da governança e da articulação
local; e as comunidades beneficiadas participam da gestão e da operação
cotidiana das infraestruturas. A definição final do modelo será consolidada
ao longo da implementação, a partir de análises técnicas, acordos
institucionais e ações socioeducativas, de mobilização comunitária e de
desenvolvimento local, visando à emancipação dos beneficiários.

\begin{fichaatividade}{Ficha-síntese da Atividade 3.2.2.1}{qua:ficha-3221}
  \campo{Código}{3.2.2.1}
  \campo{Tipo de atividade}{Estudo e mobilização}
  \campo{Forma de execução}{Contratação pública (licitação)}
  \campo{Fonte de recursos}{BIRD}
  \campo{Responsável}{IAT / IDR-Paraná / Sanepar}
  \campo{Dependências institucionais}{Prefeituras}
  \campo{Prazo estimado}{Ano 0 ao ano 6}
  \campo{Produto esperado}{Gestão dos sistemas implantada}
  \campo{Unidade de medida}{Domicílio rural}
  \campo{Evidência de comprovação}{Tarifa instituída}
  \campo{Periodicidade}{Contínua}
  \campo{Validação}{SECID}
\end{fichaatividade}

\subsubsection{Procedimentos}

A definição e a implantação do(s) modelo(s) de gestão seguirão quatro fases.

\paragraph{Fase 1 -- Definição do modelo.}
\begin{enumerate}[label=\alph*)]
  \item \textbf{critérios de seleção de áreas e comunidades}: definição dos
        critérios técnicos, sociais e ambientais para seleção das comunidades
        rurais (\autoref{sec:area});
  \item \textbf{seleção das áreas}: aplicação dos critérios para escolher as
        comunidades prioritárias;
  \item \textbf{modelo de gestão compartilhada}: definição de um modelo no
        qual Estado, municípios e comunidades compartilham responsabilidades
        pela gestão do sistema;
  \item \textbf{obrigações das partes}: definição das responsabilidades
        institucionais de cada ente envolvido.
\end{enumerate}

\paragraph{Fase 2 -- Pactuação institucional.}
\begin{enumerate}[label=\alph*), resume]
  \item \textbf{sensibilização e mobilização das prefeituras}: reuniões e
        ações institucionais para apresentar os objetivos do projeto e
        incentivar a participação dos municípios;
  \item \textbf{estruturação da governança}: criação da estrutura de
        governança do modelo, com responsabilidades, regras de funcionamento
        e mecanismos de coordenação entre os atores;
  \item \textbf{termo de cooperação}: formalização de acordo entre os
        parceiros, com os compromissos de cada instituição.
\end{enumerate}

\paragraph{Fase 3 -- Execução integrada.}
\begin{enumerate}[label=\alph*), resume]
  \item desenvolvimento do trabalho técnico-social (\atividade{3.2.2.2});
  \item execução das obras dos sistemas de abastecimento de água e de
        esgotamento sanitário (\atividade{3.2.2.3} e \atividade{3.2.2.4}).
\end{enumerate}

\paragraph{Fase 4 -- Transferência e encerramento.}
\begin{enumerate}[label=\alph*), resume]
  \item \textbf{reunião técnica de entrega}: reunião com a(s) comunidade(s) e
        com a entidade gestora para formalizar a entrega do sistema;
  \item \textbf{aceite da entidade gestora}: confirmação formal do
        recebimento da infraestrutura;
  \item \textbf{encerramento do termo de cooperação}: finalização formal do
        acordo entre os parceiros.
\end{enumerate}

\subsubsection{Etapas, produtos e evidências}

O \autoref{qua:etapas-3221} detalha as etapas de implementação dos modelos
de gestão, com os respectivos produtos, evidências e responsáveis pela
validação.

\begin{landscape}
\begin{quadro}[p]
\caption{Etapas de implementação dos modelos de gestão (\atividade{3.2.2.1})}
\label{qua:etapas-3221}
\scriptsize
\renewcommand{\arraystretch}{1.25}
\begin{tabularx}{\linewidth}{|P{2.6cm}|L|P{2.1cm}|Q{1.3cm}|P{2.5cm}|P{2cm}|P{2.6cm}|Q{1.4cm}|P{1.7cm}|}
\hline
\cab{Etapa} & \cab{Descrição técnica} & \cab{Responsável} & \cab{Prazo} &
\cab{Produto esperado} & \cab{Unidade} & \cab{Evidência} & \cab{Periodic.} &
\cab{Validação} \\ \hline
1. Estudos preliminares &
1.1 Seleção de instrumentos de gestão; 1.2 Definição de parceiros, critérios e
indicadores; 1.3 Meios de comunicação dos modelos; 1.4 Contato inicial com o
município; 1.5 Contato inicial com a comunidade; 1.6 Formalização de pré-adesão &
IAT / IDR / Prefeitura & Ano 0 a 1 & Contratos de pré-adesão &
Número de associados & Contratos assinados & Lote & IAT \\ \hline
2. Implantação &
2.1 Elaboração de documentos (estatuto, manuais, cartilhas etc.);
2.2 Enquadramento jurídico da proposta; 2.3 Adesão familiar;
2.4 Regularização da associação &
IAT / Prefeitura & Ano 1 a 2 & Associações regularizadas &
Número de associações & Registro das associações & Lote & IAT \\ \hline
3. Gestão &
3.1 Manuais de desempenho operacional; 3.2 Investimento para ação em EC;
3.3 Estrutura para monitoramento; 3.4 Treinamento &
IAT / Prefeitura / Associação & Ano 3 a 5 & Manuais, estrutura e treinamento &
Pessoas capacitadas & Lista de presença e certificados & Lote & IAT \\ \hline
4. Supervisão técnica &
4.1 Acompanhamento da implantação &
IAT / Prefeitura / Associação & Ano 2 a 6 & Supervisão e gerenciamento de obras &
Relatório de supervisão & Laudos de vistoria e testes de estanqueidade &
Contínua & IAT \\ \hline
5. Monitoramento e avaliação &
5.1 Avaliação de desempenho do processo de implementação da gestão nas
comunidades &
IAT & Ano 2 a 6 & Relatórios de conformidade técnica e ambiental &
Relatório & Dados no SIGARH / GeoPR & Semestral & UGP / MGAS \\ \hline
6. Relatórios e encerramento &
6.1 Consolidação de resultados, relatórios finais e prestação de contas &
IAT & Ano 5 e 6 & Relatório final de gerenciamento e obras concluídas &
Documento final & Termo de Recebimento Definitivo (TRD) & Única &
UGP / Banco Mundial \\ \hline
\end{tabularx}
\fonte{Adaptado de \citeonline[Quadro 28]{mop2026}.}
\end{quadro}
\end{landscape}

\pendencia{No \autoref{qua:etapas-3221}, a etapa 4 (Supervisão técnica)
repete o produto e as evidências das obras (``supervisão e gerenciamento de
obras'', ``laudos de vistoria e testes de estanqueidade''). Definir produto e
evidência próprios da implantação da gestão. A sigla ``EC'' (item 3.2) não
está definida no MOP.}

\pendencia{A ficha-síntese (Quadro 31 do MOP) adota ``tarifa instituída'' como
evidência e ``domicílio rural'' como unidade, enquanto o
\autoref{qua:etapas-3221} (Quadro 28 do MOP) usa associados, associações e
pessoas capacitadas. A validação também diverge: SECID na ficha e IAT no
quadro de etapas. Harmonizar.}

\subsubsection{Interfaces}

\begin{itemize}
  \item \atividade{3.2.2.2}: o trabalho técnico-social fornece a mobilização
        e a adesão das famílias que sustentam o modelo de gestão.
  \item \atividade{3.2.2.3} e \atividade{3.2.2.4}: os sistemas implantados
        serão transferidos à entidade gestora constituída nesta atividade.
  \item \atividade{3.2.2.6}: os modelos de gestão também se aplicam aos
        sistemas do Programa Sanepar Rural.
  \item Indicador de resultado do Componente~3 ``Associações de saneamento
        rural (RWSS) com modelos de gestão adotados''
        (\autoref{tab:metas-componente3}).
\end{itemize}
